\chapter*{Abstract}

Real-time target detection and multi-object tracking find every object of interest in a video and keep one identity for each object over time, with little delay. In tracking by detection, a detector returns scored boxes in every frame and a tracker links them into tracks. Both run with fixed settings, the operating point: the detector uses one confidence threshold, one non-maximum suppression (NMS) threshold, and one input resolution, and the tracker keeps score thresholds that were tuned for one detector. Yet scenes change within a video, above all in unmanned aerial vehicle (UAV) footage, and detectors are replaced while trackers are not retuned. Building on a published survey of real-time object detection and tracking, this thesis proposes two training-free and causal controllers that adapt these operating points online, without labels and without changing the detector or the tracker.

The first controller, Adaptive Control for Multi-Object Tracking (AC-MOT), sets the detector operating point in UAV video. Every ten frames it forms a scene complexity index from cheap image and track cues and selects one of three operating points; its values were selected by cross-validation on calibration sequences and frozen. On seven VisDrone2019-MOT sequences it raised the higher order tracking accuracy (HOTA) by 2.937 points with ByteTrack and 1.325 points with OATrack, and the multiple object tracking accuracy (MOTA) by 17.995 and 6.318 points. On UAVDT, without retuning, it raised MOTA by 17.601 and 3.267 points, at about 4.2 ms per frame.

The second controller, a self-calibrating control layer (SCL), sits between an unchanged detector and an unchanged tracker. It splits the score logits of recent frames twice with Otsu's method, keeps the tracker's scores and thresholds when the scores are clean, and otherwise chooses the candidates, rescores them, and sets the thresholds. One frozen configuration was tested on seven published trackers, four detection sources, and three datasets. It restored two trackers from 0 to about 66 HOTA under halved scores, improved both single-stage trackers on every MOT17 sequence, raised MOTA by up to 24.70 points with a noisy detector on held-out aerial data, reached higher HOTA than simple adaptive thresholds and the published ACT method, and takes 2.95 ms per frame on an ordinary processor.

The results show that the detector operating point and the score scale that the tracker receives can both be controlled online, cheaply, and without training.

\vspace{0.5cm}
\noindent\textbf{Keywords:} real-time target detection, multi-object tracking, tracking by detection, detector operating point, confidence calibration, training-free control, unmanned aerial vehicles, Otsu thresholding.
